% =============================================================================
% SECTION 11: BUSINESS CASE
% =============================================================================

\section{Business Case}

\subsection{Value Proposition}

\begin{table}[ht]
\centering
\caption{Expected Outcomes}
\label{tab:outcomes}
\begin{tabular}{lc}
\toprule
\textbf{Metric} & \textbf{Target} \\
\midrule
Takeoff time reduction & 70--80\% \\
Material waste reduction & 15--20\% $\rightarrow$ $<$5\% \\
Extraction accuracy & 95\%+ \\
Bid throughput increase & 3--5$\times$ \\
Change order reduction & 90\% \\
\bottomrule
\end{tabular}
\end{table}

\subsection{ROI Analysis}

For a mid-size contractor processing 50 projects/year:
\begin{itemize}
  \item \textbf{Time savings}: 40 hours/project $\times$ \$75/hr = \$3,000/project
  \item \textbf{Waste savings}: \$500--\$2,000/project depending on size
  \item \textbf{Annual benefit}: \$175,000--\$250,000
  \item \textbf{System cost}: \$100--\$200 per project (LLM + infrastructure)
  \item \textbf{Net annual savings}: \$165,000--\$240,000
  \item \textbf{Payback period}: $<$6 months (assuming \$50K initial setup)
\end{itemize}

\subsection{Market Context}

The U.S. construction industry represents a \$1.8 trillion annual market. The residential framing segment alone accounts for \$50+ billion in material procurement annually. Key market dynamics:
\begin{itemize}
  \item \textbf{Labor shortage}: Construction unemployment at historic lows; estimator shortage driving automation demand
  \item \textbf{Digital transformation}: Industry moving from paper plans to digital workflows
  \item \textbf{Cost pressure}: Rising material costs (lumber up 40\% since 2020) increase waste penalties
\end{itemize}

\textbf{Addressable market:}
\begin{itemize}
  \item 50,000+ general contractors in residential/light commercial
  \item 100,000+ professional estimators nationwide
  \item Average 20--50 bids per contractor per year
\end{itemize}

\subsection{Competitive Landscape}

\begin{table}[ht]
\centering
\caption{Competitive Differentiation}
\label{tab:competitors}
\begin{tabular}{lcc}
\toprule
\textbf{Feature} & \textbf{Competitors} & \textbf{Construction.AI} \\
\midrule
Automation level & Manual/partial & Full pipeline \\
Cut optimization & Separate tool & Integrated \\
Code compliance & Manual checklist & Automated + cited \\
Provenance & None & Full traceability \\
Learning & Static rules & Continuous \\
Instructions & None & Step-by-step \\
\bottomrule
\end{tabular}
\end{table}

\textbf{Competitive comparison:}
\begin{itemize}
  \item \textbf{STACK, PlanSwift}: Manual digitization with point-and-click measurement; no intelligence or optimization
  \item \textbf{Togal.AI}: AI-driven quantity extraction but no assembly reasoning, cut optimization, or build instructions
  \item \textbf{Traditional estimators}: Full-service but slow (weeks), expensive (\$2,000--\$5,000 per project), and non-scalable
  \item \textbf{Construction.AI}: End-to-end automation with KG-grounded reasoning, integrated optimization, and continuous learning
\end{itemize}

\subsection{Time to Value}

Deployment timeline for typical contractor:
\begin{itemize}
  \item \textbf{Week 1}: Initial KG setup with baseline framing rules and jurisdiction codes
  \item \textbf{Week 2--3}: Process first 3--5 projects with expert review and correction
  \item \textbf{Month 2}: System accuracy $>$90\%; begin production use with spot-checking
  \item \textbf{Month 3--6}: ROI positive as time savings compound; accuracy improves to 95\%+
\end{itemize}
